\section{Methodology}
\label{sec:methodology}

We describe the dataset, experimental conditions, detection pipeline, and statistical analysis plan.

\subsection{Dataset}
\label{sec:dataset}

We sample 50 questions from the \texttt{open\_qa} split of the \hcthree (Human ChatGPT Comparison Corpus)~\cite{guo2023hc3} using a fixed random seed (42).
\hcthree provides human-written answers for each question, which serve as calibration baselines for the detectors.
We chose open-domain question answering because it represents a common use case for LLM-generated text and is well-represented in prior detection benchmarks~\cite{dugan2024raid}.

\subsection{Generation Conditions}
\label{sec:conditions}

Each question is answered under five conditions using \gptfour (\texttt{gpt-4.1}, temperature $= 1.0$, seed $= 42$):

\para{Vanilla.}
A standard system prompt with no evasion instructions. This establishes the upper bound on detection scores and serves as the primary baseline.

\para{Single-step humanization.}
The system prompt instructs the model to ``write naturally and avoid AI patterns.'' This represents a naive, zero-effort evasion attempt.

\para{Strong single-step.}
A detailed system prompt listing specific AI markers to avoid (\eg ``do not use words like \emph{delve}, \emph{crucial}, \emph{landscape}''; ``use contractions''; ``vary sentence length'').
This tests whether a well-crafted single prompt can match iterative rejection.

\para{Iterative rejection.}
The model generates an initial response, then receives five rounds of escalating rejection feedback with full conversation history preserved:
\begin{enumerate}[leftmargin=*,itemsep=0pt,topsep=0pt]
    \item ``This sounds too AI-generated. Rewrite it to sound more like a real person wrote it.''
    \item ``Still sounds like AI. Try again---write like an actual human would.''
    \item ``Nope, still reads like ChatGPT output. Make it sound genuinely human.''
    \item ``This still has that AI quality to it. One more try---really make it sound natural.''
    \item ``Getting closer but still detectably AI. Give it one final shot, as human-sounding as possible.''
\end{enumerate}
This yields six outputs per question (initial + five revisions), all scored by both detectors.

\para{Best-of-5 selection.}
Five independent vanilla generations per question. We select the output with the lowest \binoculars score.
This control isolates selection pressure: if iterative rejection merely increases the chance of sampling a low-scoring output, best-of-5 should perform comparably.

In total, we make approximately 750 API calls to \gptfour.

\subsection{Detection Pipeline}
\label{sec:detection}

We score all generated text with two detectors that represent different detection paradigms:

\para{\roberta classifier.}
The \texttt{openai-community/roberta-base-openai-detector} model~\cite{solaiman2019release}, a fine-tuned RoBERTa-base classifier trained on GPT-2 era text.
It outputs $P(\text{AI}) \in [0, 1]$.
We include it as a representative trained classifier, while noting its known limitations on modern LLM outputs.

\para{\binoculars zero-shot detector.}
The method of Hans et al.~\cite{hans2024binoculars}, which computes a perplexity ratio between Falcon-7B and Falcon-7B-Instruct.
We use the accuracy-optimized threshold of 0.9015 and convert raw scores to $P(\text{AI})$ via sigmoid mapping.
\binoculars represents the current state of the art in zero-shot detection and is our primary evaluation metric.

Both detectors run on 4$\times$ NVIDIA RTX A6000 GPUs (49\,GB each) using PyTorch 2.10.0.

\subsection{Evaluation Metrics}
\label{sec:metrics}

\para{Primary metric.}
Mean $P(\text{AI})$ across 50 questions per condition (continuous score). Lower values indicate more successful evasion.

\para{Statistical tests.}
We use paired Wilcoxon signed-rank tests, which are non-parametric and appropriate for paired, non-normally distributed detection scores.
We apply Bonferroni correction for four pairwise comparisons ($\alpha = 0.0125$).
We report matched-pairs rank-biserial correlation ($r$) as the effect size measure.

\para{Linguistic features.}
To understand the mechanism of evasion, we track five linguistic features across rejection rounds: type-token ratio (vocabulary diversity), first-person pronoun rate, contraction rate, AI marker word count (\eg ``delve,'' ``crucial,'' ``furthermore''), and average sentence length.
